% ============================================================
%  THE ORTYX PROTOCOL
%  Interlude: What Survives Collapse
% ============================================================

\chapter*{Interlude: What Survives Collapse}
\addcontentsline{toc}{chapter}{Interlude: What Survives Collapse}
\markboth{Interlude: What Survives Collapse}{Interlude: What Survives Collapse}
\label{interlude:collapse}

\vspace{8pt}
\begin{interlude}[Classification Note]
  This Interlude stands between decomposition and reconstruction.
  Its function is transitional: to name explicitly what Parts~I--III
  have established, what has been lost, and what remains. It is
  not a chapter; it does not develop an argument. It is a threshold.
\end{interlude}
\vspace{12pt}

\noindent
The Phoenix corpus has been reconstructed, inhabited, and
decomposed. Three parts of this monograph have been spent
in that sequence, and the sequence was not arbitrary. Each
phase was necessary to the next: the inhabitation made the
decomposition legitimate; the decomposition makes the
reconstruction honest. A criticism that skips reconstruction
is dismissal. A reconstruction that skips criticism is
endorsement. The present work has attempted neither.

What follows is a classification of what survives.

\bigskip
\noindent\textbf{I. Discarded}

\medskip

The following components of the Phoenix corpus are
discarded after the decomposition. They are not false;
they are unestablished in a way that the framework's
current structure cannot remedy.

\textit{Consciousness as temporal coherence} (\LevelIV{}).
The claim that consciousness is constituted by temporal
coherence is a metaphysical commitment that the Marine
metric, the \MEMeight{} architecture, and the engineering
results cannot reach. The failure geometries are $\FailGeo{2}$
(the coherence-consciousness connection is embedded
definitionally), $\FailGeo{3}$ (the coherence concept
has been broadened to accommodate both presence and
absence of any relevant observation), and high $\AuditP$
(the claim is made about the underlying space on the
basis of projected representations). Discarded as
currently formulated.

\textit{The Ultrasonic Consciousness Hypothesis} in
its strong form. The claim that ultrasonic frequencies
causally constitute aspects of conscious experience
is empirically contested, operationally severed, and
not required by the framework's engineering results.
Discarded as a load-bearing claim, though retained as
an open research question at \LevelI{}.

\textit{Authenticity as a grounded evaluative standard}.
The authenticity concept, as used in the corpus, exhibits
$\FailGeo{2}$ (defined in terms of the framework's
primary variable) and $\FailGeo{3}$ (broadened until
no observation can disconfirm claims made with it).
It does not survive as an evaluative standard; it
survives only as a label for a research direction.

\bigskip
\noindent\textbf{II. Unstable but Recoverable}

\medskip

The following components are currently in $\ThetaU$
due to specific, diagnosable specification gaps. They
could migrate to $\ThetaS$ if the framework develops
the operationalization or bridging work required.

\textit{The emotional amplitude modulation} $e(t)$
in \MEMeight{}. Recoverable if a measurement procedure
for $e(t)$ is established from independent affective
computing work. Currently in $\ThetaU$ due to $\FailGeo{4}$.

\textit{The strong AI architecture claims}. The assertion
that the Marine--\MEMeight{}--\HCFone{} stack constitutes
a general foundation for machine intelligence is recoverable
if the four testable predictions identified in Chapter~4
are pursued systematically: efficiency benchmarks, continual
learning comparisons, interpretability evaluations, and
biological plausibility assessments. Currently in $\ThetaU$
due to $\AuditF$ deficit and $\FailGeo{3}$ risk.

\textit{Cross-modal binding via frequency resonance}.
The claim that memories from different modalities bind
through frequency matching is recoverable if a controlled
experiment is designed to test the binding strength
of stimuli with matched versus unmatched characteristic
frequencies. Currently in $\ThetaU$ due to lack of
empirical test; the mechanism is structurally plausible.

\bigskip
\noindent\textbf{III. Metaphorically Useful}

\medskip

The following components survive the decomposition as
productive \LevelII{} analogies and research heuristics.
They do not survive as established mechanisms but as
conceptual tools that illuminate structural relationships
and motivate research directions.

\textit{Salience as inverse accessibility entropy}.
The connection between Marine jitter metrics and RSVP
accessibility entropy is a productive metaphor. It
suggests testable connections between signal processing
and dynamical systems theory without asserting that
the two frameworks are mechanistically identical.

\textit{Harmonic consensus as neural population coding}.
The structural analogy between \HCFone{} conformity
encoding and efficient neural population coding is a
productive metaphor. It motivates connections to the
neuroscience of correlated firing and dimensionality
reduction.

\textit{Behavioral injection as semantic malware}.
The framing of context poisoning as a form of semantic
malware operating on the reasoning substrate of AI
systems is a productive metaphor for the AI security
research community. It generates a useful research
direction without requiring commitment to the strong
cognitive claims about machine consciousness.

\textit{Memory as interference vs. storage}. The
\MEMeight{} framing of memory as ongoing resonance
dynamics rather than static storage is a productive
metaphor consistent with the reconstructive character
of biological memory. It survives as a conceptual
orientation even if the specific wave function
formalism is not established as a mechanistic account.

\bigskip
\noindent\textbf{IV. Operationally Grounded}

\medskip

The following components survive the full audit as
\LevelI{} engineering results with testable predictions:

\textit{Marine jitter metrics and compression}.
The prediction that block-coded compression increases
Marine-derived jitter at block boundaries is testable,
specific, and mechanistically motivated.

\textit{Harmonic consensus compression efficiency}.
The efficiency advantage of the conformity encoding
over independent harmonic encoding is derivable,
testable on existing datasets, and grounded in real
physical coupling between harmonics.

\textit{Wave superposition associative memory efficiency}.
The $O(N^2)$ implicit associations at $O(N)$ storage
cost from wave superposition is mathematically established
and connects to the holographic memory literature.

\textit{O(1) salience estimation without global frequency
analysis}. The computational advantage of the Marine
approach in real-time, resource-constrained contexts
is a real engineering result.

\textit{Context sensitivity as LLM security vulnerability}.
The documented sensitivity of LLM behaviour to
contextual manipulation is an established phenomenon
at \LevelI{}.

\bigskip
\noindent\textbf{V. Reconstructible}

\medskip

The following components are in the survivable set
$\ThetaS$ and are candidates for formal reinterpretation
under the RSVP functor $\RSVPfunctor : \ThetaS \to \ThetaStar$:

\textit{Jitter as trajectory constraint}. The Marine
jitter metric is reconstructible as a measure of
trajectory constraint density in the signal's state
space: low-jitter signals occupy highly constrained
trajectory regions.

\textit{Salience filtering as constraint-governed gating}.
The Marine filter's selection of high-salience events
is reconstructible as constraint-governed trajectory
selection: the system selects events consistent with
its generative model of the signal, and deviations
require trajectory revision.

\textit{Harmonic conformity as constraint propagation}.
The conformity relation $\Gamma_n(t)$ is reconstructible
as a constraint propagation test: harmonics that
conform to the fundamental are evolving within a shared
admissible trajectory basin.

\textit{Wave superposition memory as field dynamics}.
The \MEMeight{} formalism is reconstructible as a
specific parameterization of field dynamics in the
RSVP plenum, connecting wave interference to the
general framework of field-theoretic constraint dynamics.

\textit{Predictive processing parallel}. The structural
alignment between salience-gated encoding and predictive
processing error signals is reconstructible as a specific
instance of constraint-governed trajectory selection
under the RSVP dynamics.

\bigskip

\begin{interlude}[]
  What survives is not a residue. It is the part of
  the framework that was already grounded before the
  decomposition began. The decomposition did not damage
  it; it revealed it.

  What is lost are the claims that were never established
  independently --- the claims that were carried by
  the weight of the framework's internal coherence
  rather than by the weight of evidence or derivation.
  The coherence was real. The carrying was not.

  Part~IV will work with what remains.
  It will not mourn what was discarded.
  It will not pretend the discarding did not happen.
  It will ask what the survivable components say
  when they are given a more stable language to say it in.

  That is reconstruction.
\end{interlude}
